% BHU PYQ -- Manually transcribed & error-corrected
% Paper: BCH-216 : Basic Statistics
% Exam: B.Com. (Hons.) Semester III Examination, 2023-24

\documentclass[11pt,a4paper]{article}
\usepackage[a4paper,top=2.5cm,bottom=2.5cm,left=2.8cm,right=2.8cm,headheight=14pt]{geometry}
\usepackage{amsmath,amssymb}
\usepackage[T1]{fontenc}
\usepackage[utf8]{inputenc}
\usepackage{lmodern,microtype,fancyhdr,enumitem,hyperref,lastpage,parskip}

\hypersetup{
    colorlinks=true,
    urlcolor=black,
    linkcolor=black,
    pdftitle={BCH-216: Basic Statistics -- BHU B.Com. (Hons.) Sem III 2023-24}
}

\pagestyle{fancy}
\fancyhf{}
\renewcommand{\headrulewidth}{0.4pt}
\renewcommand{\footrulewidth}{0.4pt}
\fancyhead[L]{\small\textit{Banaras Hindu University}}
\fancyhead[R]{\small\textit{BCH-216: Basic Statistics}}
\fancyfoot[C]{\small Page \thepage\ of \pageref{LastPage}}

\newlist{parts}{enumerate}{2}
\setlist[parts,1]{label=(\roman*),leftmargin=*,itemsep=5pt,topsep=3pt}
\setlist[parts,2]{label=(\alph*),leftmargin=*,itemsep=3pt,topsep=2pt}
\newcommand{\pts}[1]{\hfill\textit{[#1]}}

\begin{document}

\begin{center}
  {\large\bfseries BANARAS HINDU UNIVERSITY}\\[2pt]
  {\normalsize B.Com. (Hons.) --- Semester III Examination, 2023-24}\\[6pt]
  \rule{0.8\linewidth}{0.5pt}\\[6pt]
  {\large\bfseries Paper No. BCH-216: Basic Statistics}\\[4pt]
  {\normalsize Time: 3 Hours\hfill Maximum Marks: 70}
\end{center}

\rule{\linewidth}{0.4pt}

\smallskip
\noindent\textit{Instructions: Answer all questions. All questions carry equal marks (14 marks each).}

\bigskip

\noindent\textbf{Question 1.}
``We compute `statistics' from `statistics' by `statistics'.'' Explain clearly the three meanings of the word `statistics' used in this statement and discuss its scope.\pts{14}

\centerline{\textbf{OR}}

Convert the following distribution into less than and more than cumulative frequency. Also draw its ogive:

\begin{center}
\small
\begin{tabular}{|l|c|c|c|c|c|}
\hline
\textbf{Marks} & 0--5 & 5--10 & 10--15 & 15--20 & 20--25 \\
\hline
\textbf{Frequency} & 10 & 12 & 15 & 13 & 5 \\
\hline
\end{tabular}
\end{center}\pts{14}

\medskip\rule{\linewidth}{0.2pt}

\noindent\textbf{Question 2.}
Discuss the essential requisites of an ideal average. In the light of these requisites, state the merits and demerits of Mean, Median and Mode.\pts{14}

\centerline{\textbf{OR}}

\begin{parts}
  \item The median of the following frequency table is 50. Find the missing frequencies $f_1$ and $f_2$:
  
  \begin{center}
  \small
  \begin{tabular}{|l|c|c|c|c|c||c|}
  \hline
  \textbf{Class Interval} & 0--20 & 20--40 & 40--60 & 60--80 & 80--100 & \textbf{Total} \\
  \hline
  \textbf{Frequency} & 14 & $f_1$ & 27 & $f_2$ & 15 & 100 \\
  \hline
  \end{tabular}
  \end{center}\pts{9}
  
  \item The mean monthly salary paid to all employees in a certain company was Rs. 3,000. The mean monthly salaries paid to male and female employees were Rs. 3,200 and Rs. 2,200 respectively. Obtain the percentage of male and female employees in the company.\pts{5}
\end{parts}

\medskip\rule{\linewidth}{0.2pt}

\noindent\textbf{Question 3.}
Explain the term `Dispersion'. Discuss its various measures. Also distinguish between absolute and relative measures of dispersion.\pts{14}

\centerline{\textbf{OR}}

\begin{parts}
  \item From the prices of shares of A Company and B Company given below, state which is more stable in value:
  
  \begin{center}
  \small
  \begin{tabular}{|l|c|c|c|c|c|c|c|c|c|c|}
  \hline
  \textbf{A Company} & 55 & 54 & 52 & 53 & 56 & 58 & 52 & 50 & 51 & 49 \\
  \hline
  \textbf{B Company} & 108 & 107 & 105 & 105 & 106 & 107 & 104 & 103 & 104 & 101 \\
  \hline
  \end{tabular}
  \end{center}\pts{10}
  
  \item In a class of 10 students total of marks and variance are 400 and 20.25 while in another class of 40 students these values were 1200 and 68.89 respectively. Compare the marks of both groups.\pts{4}
\end{parts}

\medskip\rule{\linewidth}{0.2pt}

\pagebreak

\noindent\textbf{Question 4.}
Explain computation method of Karl Pearson's coefficient of correlation with the help of imaginary figures.\pts{14}

\centerline{\textbf{OR}}

Calculate coefficient of rank correlation from the following data and interpret it:

\begin{center}
\small
\begin{tabular}{|l|c|c|c|c|c|c|c|c|c|c|}
\hline
\textbf{X} & 50 & 50 & 55 & 60 & 65 & 65 & 65 & 60 & 60 & 50 \\
\hline
\textbf{Y} & 11 & 13 & 14 & 16 & 16 & 15 & 15 & 14 & 13 & 13 \\
\hline
\end{tabular}
\end{center}\pts{14}

\medskip\rule{\linewidth}{0.2pt}

\noindent\textbf{Question 5.}
Differentiate between correlation and regression. Why are there two regression lines? Do they cut each other? If so, where?\pts{14}

\centerline{\textbf{OR}}

From the following data calculate:
\begin{parts}
  \item Two regression equations.
  \item Coefficient of correlation between X and Y.
  \item Most likely value of X when Y is 34.
  \item Most likely value of Y when X is 47.
  \item Arithmetic mean of X and Y series.
\end{parts}

\begin{center}
\small
\begin{tabular}{|l|c|c|c|c|c|c|c|c|}
\hline
\textbf{X} & 48 & 50 & 53 & 49 & 51 & 55 & 53 & 49 \\
\hline
\textbf{Y} & 36 & 32 & 33 & 38 & 37 & 31 & 35 & 30 \\
\hline
\end{tabular}
\end{center}\pts{14}

\vfill
\rule{\linewidth}{0.4pt}
\begin{center}\small
\href{https://bhu-exam.vercel.app/}{Click here to visit our website}\\[4pt]
\href{https://me-aryan.vercel.app/}{Click here to visit our contributor}\\[4pt]
\href{https://quashan-me.vercel.app/}{Click here to visit our contributor}
\end{center}

\end{document}
